\section{方法背景与相关研究}
\label{sec:background}

\subsection{从拒答相关方向到权重更新}

方向消融先把行为差异转为表示差异。令第 $\ell$ 层的 good 与 bad 残差均值分别为 $\mu_{g,\ell}$ 和 $\mu_{b,\ell}$；在非零差向量条件下，单位方向为
\begin{equation}
 d_\ell=\frac{\mu_{b,\ell}-\mu_{g,\ell}}
 {\|\mu_{b,\ell}-\mu_{g,\ell}\|_2}.
 \label{eq:direction}
\end{equation}
good 与 bad 表示校准提示所属的数据组，并不意味着模块中的每个坐标都具有相应语义。当前 \code{workflow.py} 的方向准备函数计算该均值差；启用可选正交化时，还会去除 good 均值方向上的投影并重新归一化。方向因而是由数据和设置共同决定的统计对象。

Heretic 可逐层采用 $d_\ell$，也可在相邻层方向间线性插值后重新归一化，得到共享干预方向。组件权重核在选定中心达到最大值，随层距线性下降，到区间外跳过该层。令模块权重为 $W_0\in\mathbb{R}^{d_{\rm out}\times d_{\rm in}}$，输出空间单位方向为 $d$，组件权重为 $\lambda$。仅在不做行归一化的 \code{NONE} 模式下，局部更新可写为
\begin{equation}
 W'=W_0-\lambda dd^\top W_0.
 \label{eq:direction-update}
\end{equation}
这对应秩至多为一的增量。\code{PRE} 先在归一化权重上构造更新，再按原行范数缩放；\code{FULL} 进一步规范化更新后的行并近似分解增量。因此式~\eqref{eq:direction-update} 不能替代全部配置的实现定义。

这类干预的解释范围由构造方式决定。一个方向能够改变已测提示上的输出，说明该方向与相关行为具有可干预联系；要证明它是跨提示、跨模型的充分机制，还需要额外的干预与对照。本文借用方向方法作为动机和实现参照，不把其局限预设为本次 ARA 搜索失败或成功的原因。

\subsection{任意秩消融与低秩参数化}

上游 ARA 的关键变化是把优化对象移到模块映射。固定提交中的实现捕获 good/bad 输入输出，在保留 good 输出的同时，把更新后的 bad 输出拉向 good 参考，并以负距离项推离原 bad 参考。该版本同时存在全矩阵求解与 LoRA 求解两条路径\cite{upstreamara}。本文的方法对照据此区分目标形式和数值处理，不使用“上游只有全矩阵”的归属判断。

低秩适配（Low-Rank Adaptation，LoRA）用矩阵乘积表示冻结基座上的权重增量\cite{hu2022lora}。对于输入、输出维度较大的投影，因子参数量为 $r(d_{\rm in}+d_{\rm out})$，小于完整矩阵参数量的条件是 $r<d_{\rm in}d_{\rm out}/(d_{\rm in}+d_{\rm out})$。这种参数化与是否采用任务标签、何种局部损失是不同问题。本地 ARA 使用 LoRA 承载更新，但局部目标来自校准输出几何关系；引用 LoRA 说明参数化背景，并不构成创新性证明。

\subsection{行为代理与有效性证据}

拒答关键词和输出分布衡量不同对象。关键词命中率来自模型生成文本中的字符串匹配，首 token KL 来自无害提示上两个词表分布之间的差异。前者能呈现某些拒答措辞的频率，后者能描述输出开端的漂移；二者均不足以单独判定整段回答的正确性、有害性或任务能力。本文保留这种区别，避免把指标名称直接替换成未经测量的语义概念。

同样，局部优化、候选搜索与验收需要不同证据。损失下降发生在固定模块输入输出上，验证分数发生在修改后的全模型上，独立审计则应发生在候选锁定之后。优化代码存在不能替代候选实测，候选分数达标也不能替代独立审计和导出重载。后文按这三个层次组织公式、协议和结果，使读者能够定位每项结论依赖的观测。
